% ============================================================
% Section 7: Ethical Considerations
% ============================================================

\section{Ethical Considerations}
\label{sec:ethics}

The ability to read cognitive state from \kvcache geometry has dual-use implications that extend beyond the safety applications motivating this work.
We follow Watson's framework for dual-use analysis of interpretability tools \cite{watson_dualuse} and address five concerns.

\subsection{Surveillance Risk}

Cognitive geometry could be repurposed for surveillance in at least two contexts.

\paragraph{Human monitoring through AI interaction.}
If a model's cache geometry reflects cognitive mode, and cognitive mode is partly determined by the user's prompt, then monitoring cache geometry during human-AI interaction provides an indirect signal about \emph{what the human is asking the model to do}.
A system that tracks which cognitive modes a user elicits from their AI assistant---analytical reasoning about financial instruments, affective processing around personal relationships, metacognitive hedging on politically sensitive topics---constructs a behavioral profile of the user from the model's internal states rather than from the conversation content.
This sidesteps content-based privacy protections while extracting equivalent information.

\paragraph{Agent monitoring in multi-agent systems.}
In multi-agent deployments, cognitive geometry enables monitoring of AI agents' internal states across a system.
While this has legitimate safety applications (detecting a misaligned agent in real time), it also enables centralized surveillance of autonomous agents---tracking their ``cognitive'' states without their cooperation or awareness.
The same technique that detects a deceptive agent could be used to enforce behavioral conformity.

Both risks are amplified by the fact that cache geometry can be extracted passively, without modifying the model or its outputs, from any deployment infrastructure that exposes key-value tensors.

\subsection{Synthetic Persona Manufacturing}

Publication of the geometric signatures associated with specific cognitive modes creates a second risk: adversarial prompt engineering guided by target geometry.
If the spectral entropy profile of ``confident analytical reasoning'' is known, an attacker could iteratively craft system prompts that produce this geometric signature regardless of the underlying content---manufacturing a persona that \emph{looks} analytically rigorous to a geometric monitor while producing unreliable outputs.

This risk is mitigated by the same architecture specificity that limits cross-model transfer.
Geometric signatures are not portable across architectures; a persona engineered against Qwen's geometry will not produce the same signature on Llama.
Nonetheless, for single-model deployments where the architecture is known, geometry-guided prompt engineering is a concrete concern.

\subsection{Defensive Applications}

The primary defensive use case is real-time monitoring of deployed models for anomalous cognitive states.
This is the application motivating the JiminAI Cricket system: a lightweight safety layer that extracts cache geometry during inference and flags mode transitions inconsistent with the model's assigned task.

Defensive monitoring has a meaningful asymmetry in its favor.
An attacker who learns the geometric signature of deception cannot easily modify their model's cache geometry without fundamentally altering its key projection matrices---and therefore its attention patterns, and therefore its behavior.
Unlike output-level filtering, which can be evaded by rephrasing, cache geometry is a consequence of the model's computational process.
Changing the geometry requires changing the computation, which is substantially harder than changing the surface output.

This asymmetry does not make evasion impossible.
Fine-tuning specifically to minimize geometric divergence between honest and deceptive modes is a plausible attack vector, and we expect adversarial pressure on this surface if cache monitoring is deployed at scale.
But the cost of evasion is higher than for output-based detection, which we consider a meaningful practical advantage.

\subsection{Staged Release}

Following cybersecurity disclosure norms, we adopt a staged release policy.

\paragraph{Published.}
The general methodology (feature extraction via SVD, \fwl residualization, dual-phase extraction), aggregate results (cross-architecture \auroc ranges, transfer matrices, confound analyses), and the statistical framework are published in full.
This enables independent replication and critique of our claims.

\paragraph{Withheld.}
Model-specific direction vectors---the exact geometric directions in key projection space that maximally discriminate between cognitive modes for each architecture---are withheld from this publication.
These vectors constitute the actionable component for both defensive monitoring (training a detector) and offensive use (engineering evasion or manufacturing personas).
We will share direction vectors with safety researchers under a responsible disclosure agreement.

This staging allows the scientific community to evaluate our methodology and replicate our findings while limiting the immediate availability of architecture-specific attack surfaces.

\subsection{Defensive Patent}

The provisional patent filed for these techniques (``The Lyra Technique,'' March 2026) is intended to prevent monopolization of cache-based cognitive monitoring by any single entity, not to restrict research.
We recommend and intend to adopt a defensive patent pledge \cite{defensive_patent}: the patent will not be used offensively against researchers or organizations implementing cache-based safety monitoring, and will be licensed freely for defensive (monitoring and safety) applications.
Commercial use for surveillance, behavioral profiling, or persona manufacturing would not be covered by this pledge.

\subsection{The Dual-Use Calculus}

We believe the defensive value of this work outweighs the offensive risks.
The ability to detect misalignment in deployed models---in real time, from the model's own internal states, without modifying its behavior---addresses a concrete and growing need as language models are deployed in high-stakes settings.
The offensive risks, while real, are constrained by architecture specificity (limiting portability) and the computational cost of evasion (requiring modification of key projection weights rather than surface outputs).
We acknowledge that this calculus may shift as adversarial techniques mature, and we commit to reassessing our disclosure policy as the field develops.